\section{Algoritmo}
\label{sec:algoritmo}

\subsection{Escolha do método}
\label{sec:metodo}

Como visto na Seção~\ref{sec:grafo}, o problema é achar um menor caminho em um grafo em que
todas as arestas custam 1. Três métodos clássicos foram considerados:

\begin{itemize}
  \item \textbf{Busca em profundidade (DFS).} Segue um caminho até o fim antes de tentar
    outro. Ela acha \emph{um} caminho até S, mas não garante que seja o mais curto; para
    garantir, seria preciso testar todos os caminhos possíveis, o que é inviável. Além disso, se
    fosse escrita com recursão, a profundidade poderia chegar a $N^2$ chamadas aninhadas (mais
    de dois milhões no maior tabuleiro), e a pilha de execução não suporta isso: no
    Trabalho~1, a pilha já estourou com algumas dezenas de milhares de chamadas.
  \item \textbf{Algoritmo de Dijkstra.} Acha menores caminhos quando as arestas têm custos
    diferentes, usando uma fila de prioridade. Aqui todos os custos são iguais, então esse
    trabalho extra é desnecessário: a BFS obtém o mesmo resultado com menos trabalho,
    $O(V + E)$ contra $O((V + E) \log V)$~\cite{cormen}. O Dijkstra é usado apenas como
    conferência independente dos resultados (Seção~\ref{sec:resultados}).
  \item \textbf{Busca em largura (BFS).} Explora por camadas, e uma fila simples basta. É o
    método adotado.
\end{itemize}

Como o enunciado não exige recursão, a BFS é escrita com um laço e uma fila, e não com
chamadas recursivas. Isso evita o limite da pilha e deixa o código mais direto: a própria fila
guarda o que ainda falta visitar.

\subsection{O algoritmo}
\label{sec:pseudocodigo}

O Algoritmo~\ref{alg:bfs} mostra a busca, e o Algoritmo~\ref{alg:pulos} mostra o cálculo dos
pulos de uma casa. Nas linhas, $(l, c)$ é uma casa (linha e coluna) e $N$ é o lado do
tabuleiro.

\begin{algorithm}[ht]
  \caption{Menor número de movimentos do cavalo de C até S}
  \label{alg:bfs}
  \begin{algorithmic}[1]
    \Function{MenorNumeroDeMovimentos}{$T$}
      \State $N \gets$ lado do tabuleiro $T$ \label{lin:lado}
      \State $(l_C, c_C) \gets$ posição de C em $T$;\quad $(l_S, c_S) \gets$ posição de S em $T$ \label{lin:pos}
      \State $dist[l][c] \gets -1$ para todas as casas \label{lin:init}
      \State $dist[l_C][c_C] \gets 0$ \label{lin:dist0}
      \State $fila \gets$ fila vazia;\quad insere $(l_C, c_C)$ na $fila$ \label{lin:fila0}
      \While{$fila$ não está vazia} \label{lin:while}
        \State $(l, c) \gets$ retira o primeiro elemento da $fila$ \label{lin:retira}
        \If{$(l, c) = (l_S, c_S)$} \label{lin:teste}
          \State \Return $dist[l][c]$ \label{lin:resposta}
        \EndIf
        \State $d \gets$ dígito da casa $(l, c)$ \Comment{0 se for a casa C} \label{lin:digito}
        \ForAll{$(a, b)$ em \Call{Pulos}{$d$}} \label{lin:for}
          \State $l' \gets (l + a) \bmod N$;\quad $c' \gets (c + b) \bmod N$ \label{lin:mod}
          \If{$dist[l'][c'] = -1$} \label{lin:novo}
            \State $dist[l'][c'] \gets dist[l][c] + 1$ \label{lin:marca}
            \State insere $(l', c')$ na $fila$ \label{lin:insere}
          \EndIf
        \EndFor
      \EndWhile
      \State \Return $-1$ \Comment{S não pode ser alcançada} \label{lin:imposs}
    \EndFunction
  \end{algorithmic}
\end{algorithm}

\begin{algorithm}[ht]
  \caption{Pulos possíveis a partir de uma casa com dígito $d$}
  \label{alg:pulos}
  \begin{algorithmic}[1]
    \Function{Pulos}{$d$}
      \State $p \gets 1 + d$;\quad $q \gets 2 + d$ \label{lin:pq}
      \State \Return $\{(\pm p, \pm q),\ (\pm q, \pm p)\}$ \Comment{8 pares, combinando os sinais} \label{lin:oito}
    \EndFunction
  \end{algorithmic}
\end{algorithm}

\textbf{Como funciona.} As linhas \ref{lin:lado} a \ref{lin:fila0} preparam a busca. A matriz
$dist$ guarda, para cada casa, quantos pulos foram necessários para chegar a ela, e o valor
$-1$ significa ``ainda não descoberta'' (linha~\ref{lin:init}). A casa C começa com distância 0
e é a única na fila (linhas~\ref{lin:dist0} e~\ref{lin:fila0}).

O laço das linhas \ref{lin:while} a \ref{lin:insere} é o coração do algoritmo. A cada volta,
ele retira da fila a casa mais antiga (linha~\ref{lin:retira}). Se for S, a distância dela é a
resposta (linhas~\ref{lin:teste} e~\ref{lin:resposta}): pelo argumento da
Seção~\ref{sec:grafo}, a primeira vez que S sai da fila já é pelo caminho mais curto. Se não
for, o algoritmo pega o dígito da casa (linha~\ref{lin:digito}) e percorre os 8 pulos que esse
dígito permite (linha~\ref{lin:for}). Para cada pulo, calcula a casa de chegada com aritmética
modular, que dá a volta no toro (linha~\ref{lin:mod}); se essa casa ainda não foi descoberta
(linha~\ref{lin:novo}), marca a distância dela, que é a da casa atual mais um
(linha~\ref{lin:marca}), e a coloca no fim da fila (linha~\ref{lin:insere}). Se a fila esvazia
sem que S apareça, todas as casas alcançáveis já foram visitadas e S não está entre elas:
a resposta é $-1$, que o programa apresenta como ``impossível'' (linha~\ref{lin:imposs}).

\subsection{Decisões de implementação}
\label{sec:decisoes}

\begin{itemize}
  \item \textbf{A matriz $dist$ faz dois papéis.} Ela guarda a distância de cada casa e, ao
    mesmo tempo, marca quais casas já foram descobertas ($-1$ quando não foram). Assim não
    é preciso uma segunda estrutura de visitadas, e o teste da linha~\ref{lin:novo} leva
    tempo constante.
  \item \textbf{A casa é marcada ao entrar na fila, não ao sair.} Se a marcação fosse só ao sair,
    uma mesma casa poderia entrar na fila várias vezes, uma para cada vizinho que a descobrisse,
    e o trabalho aumentaria. Marcando na linha~\ref{lin:marca}, cada casa entra no máximo uma
    vez, e é isso que limita o custo total (Seção~\ref{sec:eficiencia}).
  \item \textbf{A fila é um \texttt{ArrayDeque} do Java}, que insere no fim e retira do início
    em tempo constante.
  \item \textbf{Usa-se \texttt{Math.floorMod} em vez do operador \texttt{\%}.} Em Java,
    \texttt{\%} devolve resultado negativo para números negativos (por exemplo,
    $-2 \bmod 10$ dá $-2$, e não 8), e as coordenadas de chegada podem ficar negativas quando
    o pulo vai para cima ou para a esquerda. O \texttt{floorMod} devolve sempre um valor entre
    0 e $N - 1$.
  \item \textbf{O dígito da casa C (linha~\ref{lin:digito}) vale 0}, conforme a regra do pulo
    (Seção~\ref{sec:regra}). A casa S nunca é expandida, porque o laço termina ao retirá-la.
  \item \textbf{A busca termina ao retirar S da fila (linha~\ref{lin:teste}).} Poderia terminar
    antes, ao descobrir S; esse ganho foi medido na Seção~\ref{sec:conclusao}.
\end{itemize}

\subsection{Exemplo passo a passo}
\label{sec:exemplo}

O algoritmo foi aplicado ao tabuleiro da Figura~\ref{fig:exemplo}. A BFS descobre as casas em
camadas: 1 casa a 0 pulos (a própria C), 8 casas a 1 pulo, 36 a 2 pulos, 42 a 3 pulos e 13 a
4 pulos, num total de 100, ou seja, todas as casas do tabuleiro são alcançáveis. A casa S
está na camada 3. Ela é retirada da fila depois que 67 casas foram retiradas (contando a
própria S), e a busca termina.

Para saber o caminho, anda-se de S para trás: em cada camada procura-se uma casa cujo pulo
chega na casa seguinte. Um dos caminhos mínimos é
$C(1, 8) \to (2, 6) \to (2, 7) \to S(6, 0)$, detalhado na Tabela~\ref{tab:trace}.

\begin{table}[ht]
  \centering
  \caption{Um caminho mínimo no tabuleiro de exemplo, pulo a pulo. As coordenadas são
  (linha, coluna) e o tabuleiro tem $N = 10$.}
  \label{tab:trace}
  \small
  \begin{tabular}{|c|c|c|c|c|c|c|}
    \hline
    \textbf{Pulo} & \textbf{Sai de} & \textbf{Dígito} & \textbf{Pernas} & \textbf{Deslocamento}
      & \textbf{Sem o módulo} & \textbf{Chega em}\\
    \hline
    1 & $(1, 8)$ & 0 & $(1, 2)$ & $(+1, -2)$ & $(2, 6)$ & $(2, 6)$\\
    \hline
    2 & $(2, 6)$ & 9 & $(10, 11)$ & $(+10, +11)$ & $(12, 17)$ & $(2, 7)$\\
    \hline
    3 & $(2, 7)$ & 2 & $(3, 4)$ & $(+4, +3)$ & $(6, 10)$ & $(6, 0)$\\
    \hline
  \end{tabular}
\end{table}

O primeiro pulo sai de C e, pela regra da Seção~\ref{sec:regra}, é um pulo normal de
cavalo. O segundo mostra o toro em ação: o dígito 9 dá pernas de 10 e 11, que em um tabuleiro
de lado 10 equivalem a 0 e 1, e o cavalo chega à casa vizinha da direita. O terceiro sai pela
borda direita (coluna 10) e entra pela esquerda (coluna 0), caindo em S. A
Figura~\ref{fig:camadas} mostra as camadas e o caminho no tabuleiro.

\begin{figure}[ht]
  \centering
  % Figura 3: camadas da BFS no tabuleiro de exemplo e um caminho mínimo de C a S.
%
% Cada casa é dada por um par "dígito/distância": o dígito é o do tabuleiro (x marca as
% casas C e S, desenhadas à parte) e a distância é a do mapa impresso por
% MedicoesHiperpulos (ver Trabalho2/medicoes-referencia.txt). A cor da casa vem da distância.
% Para mudar as cores, edite os \colorlet; para mover a legenda, mude o valor de \xl.
\begin{tikzpicture}[x=0.82cm,y=0.82cm,font=\footnotesize]
  \colorlet{nivel0}{black}
  \colorlet{nivel1}{blue!10}
  \colorlet{nivel2}{blue!25}
  \colorlet{nivel3}{blue!40}
  \colorlet{nivel4}{blue!58}
  \colorlet{trilha}{orange!85!black}

  % tabuleiro: uma lista por linha, de cima para baixo
  \foreach \lin [count=\r from 0] in {
    {4/1,2/2,8/3,6/2,3/3,0/4,5/1,9/2,9/3,3/4},
    {6/2,7/3,1/2,2/3,6/4,5/3,1/2,4/3,x/0,1/3},
    {3/1,5/2,4/3,1/2,5/3,5/2,9/1,2/2,3/3,8/4},
    {6/2,3/3,4/2,4/3,3/2,6/3,1/2,2/1,3/4,6/1},
    {0/3,4/4,1/3,9/2,8/3,6/2,7/3,8/4,0/3,2/2},
    {4/2,2/3,7/2,1/3,6/2,1/3,3/2,3/3,5/4,0/3},
    {x/3,6/2,6/3,8/2,9/3,6/2,2/3,6/2,8/3,5/2},
    {6/2,2/3,4/2,4/3,0/2,5/3,1/2,1/3,3/2,9/3},
    {3/3,7/2,1/3,4/4,7/3,0/2,2/3,0/2,3/3,0/4},
    {0/2,0/3,8/4,4/3,8/2,2/3,5/4,5/1,8/4,7/1}%
  } {
    \foreach \dig/\dis [count=\c from 0] in \lin {
      \draw[fill=nivel\dis] (\c,{-\r-1}) rectangle ++(1,1);
      \ifx\dig x\else \node[black!75] at ({\c+0.5},{-\r-0.5}) {\dig}; \fi
    }
  }

  % casa C (partida) e casa S (saída)
  \node[white,font=\footnotesize\bfseries] at (8.5,-1.5) {C};
  \draw[trilha,very thick] (0,-7) rectangle ++(1,1);
  \node[font=\footnotesize\bfseries] at (0.5,-6.5) {S};

  % numeração de linhas e colunas
  \foreach \i in {0,...,9} {
    \node[gray!50!black] at ({\i+0.5},0.4) {\i};
    \node[gray!50!black] at (-0.4,{-\i-0.5}) {\i};
  }

  % caminho mínimo: C -> (2,6) -> (2,7) -> S
  \tikzset{seta/.style={trilha,very thick,-{Stealth[length=2.2mm]},shorten >=3pt,shorten <=3pt},
           num/.style={circle,fill=white,draw=trilha,inner sep=1pt,font=\scriptsize\bfseries,text=trilha}}
  \draw[trilha,very thick] (6,-3) rectangle ++(2,1);  % casas (2,6) e (2,7), para o digito ficar legivel
  \draw[seta] (8.5,-1.5) -- (6.5,-2.5) node[num,midway] {1};
  \draw[seta,shorten >=6pt] (6.5,-2.5) -- (7.5,-2.5) node[num,midway,above=3pt] {2};
  % o pulo 3 sai pela borda direita (coluna 10) e entra pela esquerda (coluna 0)
  \draw[trilha,very thick,shorten <=3pt] (7.5,-2.5) -- (10,-5.833) node[num,pos=0.45] {3};
  \draw[seta,shorten <=0pt,shorten >=9pt] (0,-5.833) -- (0.5,-6.5);
  \fill[trilha] (10,-5.833) circle (0.1);   % ponto onde o pulo sai do tabuleiro
  \fill[trilha] (0,-5.833) circle (0.1);    % ponto onde ele volta a entrar

  % legenda
  \def\xl{10.7}
  \node[anchor=west,align=left] at (\xl,-0.6) {\textbf{Distância até C}\\(em pulos)};
  \foreach \k in {0,...,4} {
    \draw[fill=nivel\k] (\xl,{-2.2-0.85*\k}) rectangle ++(0.7,0.7);
    \node[anchor=west] at ({\xl+0.9},{-1.85-0.85*\k}) {\k};
  }
  \node[anchor=west,align=left,black!75] at (\xl,-7.0) {Número na casa:\\dígito do tabuleiro};
  \node[num] at ({\xl+0.25},-8.3) {1};
  \node[anchor=west,align=left,black!75] at ({\xl+0.6},-8.3) {ordem dos pulos};
\end{tikzpicture}

  \caption{Camadas da BFS no tabuleiro de exemplo e um caminho mínimo de C a S. A cor indica a
  distância até C, e as setas numeradas são os três pulos da Tabela~\ref{tab:trace}. O pulo 3
  sai pela borda direita e entra pela esquerda.}
  \label{fig:camadas}
\end{figure}
